\section{前端性能优化工具}\label{ux524dux7aefux6027ux80fdux4f18ux5316ux5de5ux5177}

前端性能优化是智慧水利平台用户体验的关键因素，良好的性能能够提高用户满意度、降低运营成本并提升系统可靠性。本文详细介绍智慧水利平台前端性能优化工具及最佳实践。

\subsection{1. 性能分析工具}\label{ux6027ux80fdux5206ux6790ux5de5ux5177}

\subsubsection{1.1
浏览器开发者工具}\label{ux6d4fux89c8ux5668ux5f00ux53d1ux8005ux5de5ux5177}

主流浏览器内置的开发者工具是进行性能分析的首选工具。

\paragraph{Chrome DevTools
Performance面板}\label{chrome-devtools-performanceux9762ux677f}

Chrome DevTools的Performance面板提供全面的性能分析能力：

\begin{enumerate}
\def\labelenumi{\arabic{enumi}.}
\tightlist
\item
  \textbf{性能记录与分析}

  \begin{itemize}
  \tightlist
  \item
    使用方法：打开DevTools \textgreater{} Performance \textgreater{}
    点击录制按钮
  \item
    功能：捕获页面加载和用户交互期间的性能数据
  \end{itemize}
\item
  \textbf{主要指标}

  \begin{itemize}
  \tightlist
  \item
    FPS（帧率）：页面动画流畅度
  \item
    CPU使用率：按类别划分的CPU活动
  \item
    网络请求时间线：资源加载时间
  \item
    主线程活动：JavaScript执行、样式计算、布局等
  \item
    内存使用情况：内存占用变化
  \end{itemize}
\item
  \textbf{智慧水利平台优化案例}
\end{enumerate}

\begin{lstlisting}
// 使用performance.mark()和performance.measure()进行自定义性能标记
// 在水位图表初始化前
performance.mark('waterLevelChart-start');

// 初始化水位图表
initWaterLevelChart(data);

// 在水位图表初始化后
performance.mark('waterLevelChart-end');

// 计算图表初始化耗时
performance.measure(
  'waterLevelChart-init', 
  'waterLevelChart-start', 
  'waterLevelChart-end'
);

// 在DevTools中查看这些自定义标记
\end{lstlisting}

\paragraph{Lighthouse}\label{lighthouse}

Lighthouse是Google开发的自动化性能审计工具：

\begin{enumerate}
\def\labelenumi{\arabic{enumi}.}
\tightlist
\item
  \textbf{主要性能指标}

  \begin{itemize}
  \tightlist
  \item
    First Contentful Paint (FCP)：首次内容绘制
  \item
    Largest Contentful Paint (LCP)：最大内容绘制
  \item
    Cumulative Layout Shift (CLS)：累积布局偏移
  \item
    Total Blocking Time (TBT)：总阻塞时间
  \item
    Time to Interactive (TTI)：可交互时间
  \end{itemize}
\item
  \textbf{使用方法}

  \begin{itemize}
  \tightlist
  \item
    Chrome DevTools \textgreater{} Lighthouse标签
  \item
    命令行工具：\passthrough{\lstinline!npx lighthouse https://your-site.com!}
  \item
    CI/CD集成：使用lighthouse-ci
  \end{itemize}
\item
  \textbf{智慧水利平台优化参考}

  \begin{itemize}
  \tightlist
  \item
    LCP \textless{} 2.5秒：首屏大图表快速加载
  \item
    CLS \textless{} 0.1：确保数据加载不引起布局跳动
  \item
    TTI \textless{} 3.8秒：确保大屏幕监控快速可交互
  \end{itemize}
\end{enumerate}

\subsubsection{1.2 Web
Vitals与性能监控}\label{web-vitalsux4e0eux6027ux80fdux76d1ux63a7}

Web Vitals是Google提出的核心Web性能指标，对用户体验至关重要。

\paragraph{集成Web Vitals监控}\label{ux96c6ux6210web-vitalsux76d1ux63a7}

\begin{lstlisting}
// web-vitals.js
import { getLCP, getFID, getCLS } from 'web-vitals';

function sendToAnalytics(metric) {
  const body = JSON.stringify({
    name: metric.name,
    value: metric.value,
    id: metric.id,
    page: location.pathname,
    platform: 'web'
  });
  
  // 发送到分析服务
  navigator.sendBeacon('/analytics', body);
}

// 监测核心Web指标
getCLS(sendToAnalytics);
getFID(sendToAnalytics);
getLCP(sendToAnalytics);
\end{lstlisting}

\paragraph{性能监控系统}\label{ux6027ux80fdux76d1ux63a7ux7cfbux7edf}

\begin{enumerate}
\def\labelenumi{\arabic{enumi}.}
\tightlist
\item
  \textbf{自定义性能监控}
\end{enumerate}

\begin{lstlisting}
// performance-monitoring.js
class PerformanceMonitor {
  constructor() {
    this.metrics = {};
    this.initObservers();
  }
  
  initObservers() {
    // 页面加载性能
    window.addEventListener('load', () => {
      setTimeout(() => {
        const navigation = performance.getEntriesByType('navigation')[0];
        const paintEntries = performance.getEntriesByType('paint');
        
        this.metrics.domComplete = navigation.domComplete;
        this.metrics.loadEventEnd = navigation.loadEventEnd;
        this.metrics.domInteractive = navigation.domInteractive;
        
        const fcp = paintEntries.find(entry => entry.name === 'first-contentful-paint');
        if (fcp) this.metrics.fcp = fcp.startTime;
        
        this.sendMetrics();
      }, 0);
    });
    
    // 资源加载性能
    const resourceObserver = new PerformanceObserver((list) => {
      const resources = list.getEntries().filter(entry => 
        entry.initiatorType !== 'fetch' && 
        entry.initiatorType !== 'xmlhttprequest'
      );
      
      this.metrics.resources = resources.map(r => ({
        name: r.name,
        type: r.initiatorType,
        size: r.transferSize,
        duration: r.duration
      }));
    });
    
    resourceObserver.observe({ entryTypes: ['resource'] });
  }
  
  sendMetrics() {
    fetch('/api/performance', {
      method: 'POST',
      headers: { 'Content-Type': 'application/json' },
      body: JSON.stringify(this.metrics)
    });
  }
}

// 初始化监控
new PerformanceMonitor();
\end{lstlisting}

\begin{enumerate}
\def\labelenumi{\arabic{enumi}.}
\setcounter{enumi}{1}
\tightlist
\item
  \textbf{第三方监控服务集成}

  \begin{itemize}
  \tightlist
  \item
    New Relic
  \item
    Datadog RUM
  \item
    Sentry Performance
  \end{itemize}
\end{enumerate}

\subsubsection{1.3 WebPageTest}\label{webpagetest}

WebPageTest是全面的性能测试工具，可从不同地点和网络条件测试网站性能。

\begin{enumerate}
\def\labelenumi{\arabic{enumi}.}
\tightlist
\item
  \textbf{主要功能}

  \begin{itemize}
  \tightlist
  \item
    多地点测试
  \item
    多设备测试
  \item
    网络节流模拟
  \item
    视频捕获和视觉比较
  \item
    瀑布图分析
  \end{itemize}
\item
  \textbf{关键性能指标}

  \begin{itemize}
  \tightlist
  \item
    Time to First Byte (TTFB)
  \item
    First Paint
  \item
    Document Complete
  \item
    Fully Loaded
  \item
    Speed Index
  \end{itemize}
\item
  \textbf{内容分发优化}

  \begin{itemize}
  \tightlist
  \item
    CDN使用分析
  \item
    缓存策略评估
  \item
    HTTP/2和HTTP/3支持检测
  \end{itemize}
\end{enumerate}

\subsection{2.
监控与错误追踪}\label{ux76d1ux63a7ux4e0eux9519ux8befux8ffdux8e2a}

\subsubsection{2.1 Sentry}\label{sentry}

Sentry是应用程序监控平台，可用于错误追踪和性能监控。

\paragraph{Sentry集成配置}\label{sentryux96c6ux6210ux914dux7f6e}

\begin{lstlisting}
// 初始化Sentry
import * as Sentry from '@sentry/vue';
import { BrowserTracing } from '@sentry/tracing';
import Vue from 'vue';
import router from './router';

Sentry.init({
  Vue,
  dsn: 'https://examplePublicKey@o0.ingest.sentry.io/0',
  
  integrations: [
    new BrowserTracing({
      routingInstrumentation: Sentry.vueRouterInstrumentation(router),
      tracingOrigins: ['localhost', 'water-platform.example.com']
    }),
  ],
  
  // 跟踪样本率
  tracesSampleRate: 0.2,
  
  // 环境信息
  environment: process.env.NODE_ENV,
  
  // 发布版本
  release: 'water-platform@' + process.env.VUE_APP_VERSION,
  
  // 忽略特定错误
  ignoreErrors: [
    'Network Error',
    /ChunkLoadError/,
    /loading chunk/i
  ],
  
  // 添加水利平台特定标签
  beforeSend(event) {
    event.tags = {
      ...event.tags,
      module: getCurrentModuleName()
    };
    return event;
  }
});

// 添加用户上下文信息
if (userStore.isLoggedIn) {
  Sentry.setUser({
    id: userStore.userId,
    username: userStore.username,
    role: userStore.role
  });
}

// 添加自定义事务
export function monitorOperation(operationName, callback) {
  const transaction = Sentry.startTransaction({
    name: operationName
  });
  
  Sentry.configureScope(scope => {
    scope.setSpan(transaction);
  });
  
  try {
    return callback();
  } catch (error) {
    Sentry.captureException(error);
    throw error;
  } finally {
    transaction.finish();
  }
}
\end{lstlisting}

\paragraph{自定义错误边界组件}\label{ux81eaux5b9aux4e49ux9519ux8befux8fb9ux754cux7ec4ux4ef6}

\begin{lstlisting}
<!-- ErrorBoundary.vue -->
<template>
  <div>
    <slot v-if="!error"></slot>
    <div v-else class="error-container">
      <h3>组件加载失败</h3>
      <p>{{ errorMessage }}</p>
      <button @click="resetError">重试</button>
    </div>
  </div>
</template>

<script>
import * as Sentry from '@sentry/vue';

export default {
  name: 'ErrorBoundary',
  props: {
    componentName: {
      type: String,
      default: '未知组件'
    }
  },
  data() {
    return {
      error: null,
      errorMessage: ''
    };
  },
  errorCaptured(err, vm, info) {
    this.error = err;
    this.errorMessage = this.formatErrorMessage(err);
    
    // 向Sentry报告错误
    Sentry.captureException(err, {
      tags: {
        componentName: this.componentName,
        info
      }
    });
    
    // 阻止错误继续传播
    return false;
  },
  methods: {
    resetError() {
      this.error = null;
      this.errorMessage = '';
    },
    formatErrorMessage(error) {
      if (process.env.NODE_ENV === 'production') {
        return '应用程序发生错误，已记录并将尽快修复。';
      }
      return error.message || String(error);
    }
  }
};
</script>
\end{lstlisting}

\subsubsection{2.2 Google Analytics}\label{google-analytics}

Google Analytics可用于跟踪用户行为和性能指标。

\paragraph{智慧水利平台GA配置}\label{ux667aux6167ux6c34ux5229ux5e73ux53f0gaux914dux7f6e}

\begin{lstlisting}
// analytics.js
import { reactive } from 'vue';

// GA初始化代码
function initGA() {
  const script = document.createElement('script');
  script.async = true;
  script.src = `https://www.googletagmanager.com/gtag/js?id=${process.env.VUE_APP_GA_ID}`;
  document.head.appendChild(script);
  
  window.dataLayer = window.dataLayer || [];
  function gtag() {
    dataLayer.push(arguments);
  }
  gtag('js', new Date());
  gtag('config', process.env.VUE_APP_GA_ID, {
    send_page_view: false,
    custom_map: {
      dimension1: 'user_role',
      dimension2: 'region_code',
      metric1: 'response_time'
    }
  });
  
  return gtag;
}

const gtag = initGA();

// 创建分析服务
export const analytics = reactive({
  // 页面访问
  pageView(path, title) {
    gtag('event', 'page_view', {
      page_path: path,
      page_title: title
    });
  },
  
  // 事件追踪
  trackEvent(category, action, label, value) {
    gtag('event', action, {
      event_category: category,
      event_label: label,
      value: value
    });
  },
  
  // 用户设置
  setUserProperties(properties) {
    gtag('set', 'user_properties', properties);
  },
  
  // 性能指标跟踪
  trackPerformance(metrics) {
    gtag('event', 'performance', metrics);
  },
  
  // 异常追踪
  trackException(description, fatal = false) {
    gtag('event', 'exception', {
      description,
      fatal
    });
  }
});

// 与Vue Router集成
export function setupAnalyticsRouter(router) {
  router.afterEach((to) => {
    analytics.pageView(to.fullPath, to.meta.title || to.name);
  });
}
\end{lstlisting}

\subsubsection{2.3
自定义性能指标}\label{ux81eaux5b9aux4e49ux6027ux80fdux6307ux6807}

智慧水利平台特定业务场景的性能指标。

\paragraph{自定义性能指标收集}\label{ux81eaux5b9aux4e49ux6027ux80fdux6307ux6807ux6536ux96c6}

\begin{lstlisting}
// performance-metrics.js
class WaterPlatformMetrics {
  constructor() {
    this.metrics = {};
    this.timers = {};
  }
  
  // 开始计时
  startTimer(key) {
    this.timers[key] = performance.now();
  }
  
  // 结束计时并记录指标
  endTimer(key) {
    if (!this.timers[key]) return;
    const duration = performance.now() - this.timers[key];
    this.metrics[key] = duration;
    delete this.timers[key];
    return duration;
  }
  
  // 记录水位数据加载时间
  trackWaterLevelDataLoad(stationCount, duration) {
    this.metrics.waterLevelLoad = {
      stations: stationCount,
      duration,
      timestamp: Date.now()
    };
  }
  
  // 记录图表渲染时间
  trackChartRender(chartType, dataPoints, duration) {
    if (!this.metrics.chartRenders) {
      this.metrics.chartRenders = [];
    }
    
    this.metrics.chartRenders.push({
      type: chartType,
      dataPoints,
      duration,
      timestamp: Date.now()
    });
  }
  
  // 记录用户交互响应时间
  trackInteraction(actionType, duration) {
    if (!this.metrics.interactions) {
      this.metrics.interactions = [];
    }
    
    this.metrics.interactions.push({
      type: actionType,
      duration,
      timestamp: Date.now()
    });
  }
  
  // 获取所有指标
  getAllMetrics() {
    return this.metrics;
  }
  
  // 发送指标到服务器
  sendMetrics() {
    fetch('/api/client-metrics', {
      method: 'POST',
      headers: { 'Content-Type': 'application/json' },
      body: JSON.stringify(this.getAllMetrics())
    });
  }
  
  // 定期发送指标
  startAutoCollection(intervalMs = 60000) {
    setInterval(() => this.sendMetrics(), intervalMs);
  }
}

// 导出单例
export const waterMetrics = new WaterPlatformMetrics();
\end{lstlisting}

\subsection{3.
打包分析与优化}\label{ux6253ux5305ux5206ux6790ux4e0eux4f18ux5316}

\subsubsection{3.1
webpack-bundle-analyzer}\label{webpack-bundle-analyzer}

webpack-bundle-analyzer可视化分析Webpack包内容和大小。

\paragraph{配置与使用}\label{ux914dux7f6eux4e0eux4f7fux7528}

\begin{lstlisting}
// vue.config.js (Vue CLI项目)
const { BundleAnalyzerPlugin } = require('webpack-bundle-analyzer');

module.exports = {
  chainWebpack: config => {
    if (process.env.ANALYZE) {
      config.plugin('webpack-bundle-analyzer')
        .use(BundleAnalyzerPlugin, [{
          analyzerMode: 'static',
          reportFilename: 'bundle-report.html',
          openAnalyzer: false
        }]);
    }
  }
};
\end{lstlisting}

运行命令：\passthrough{\lstinline!ANALYZE=true npm run build!}

\paragraph{优化实践}\label{ux4f18ux5316ux5b9eux8df5}

基于分析结果的优化策略：

\begin{enumerate}
\def\labelenumi{\arabic{enumi}.}
\item
  \textbf{拆分大型库}

\begin{lstlisting}
// vue.config.js
module.exports = {
  configureWebpack: {
    optimization: {
      splitChunks: {
        cacheGroups: {
          echarts: {
            name: 'chunk-echarts',
            test: /[\\/]node_modules[\\/]echarts[\\/]/,
            priority: 20,
            chunks: 'all'
          },
          elementUI: {
            name: 'chunk-elementui',
            test: /[\\/]node_modules[\\/]element-ui[\\/]/,
            priority: 20,
            chunks: 'all'
          }
        }
      }
    }
  }
};
\end{lstlisting}
\item
  \textbf{懒加载组件}

\begin{lstlisting}
// router/index.js
const routes = [
  {
    path: '/water-analysis',
    name: 'WaterAnalysis',
    component: () => import(/* webpackChunkName: "water-analysis" */ '../views/WaterAnalysis.vue')
  },
  {
    path: '/monitoring-dashboard',
    name: 'MonitoringDashboard',
    component: () => import(/* webpackChunkName: "monitoring" */ '../views/MonitoringDashboard.vue')
  }
];
\end{lstlisting}
\item
  \textbf{树摇动优化}

\begin{lstlisting}
// 按需导入
import { LineChart, BarChart } from 'echarts/charts';
import { 
  GridComponent,
  TooltipComponent,
  LegendComponent
} from 'echarts/components';

// 而不是
// import * as echarts from 'echarts';
\end{lstlisting}
\end{enumerate}

\subsubsection{3.2 Source Map Explorer}\label{source-map-explorer}

Source Map Explorer利用源映射分析打包后的JavaScript文件。

\paragraph{使用方法}\label{ux4f7fux7528ux65b9ux6cd5}

\begin{lstlisting}[language=bash]
# 安装
npm install --save-dev source-map-explorer

# 生成带sourceMap的构建
GENERATE_SOURCEMAP=true npm run build

# 分析特定文件
npx source-map-explorer dist/js/chunk-vendors.js
\end{lstlisting}

\paragraph{配置示例}\label{ux914dux7f6eux793aux4f8b}

\begin{lstlisting}
// package.json
{
  "scripts": {
    "analyze": "source-map-explorer 'dist/js/*.js'",
    "build:analyze": "GENERATE_SOURCEMAP=true npm run build && npm run analyze"
  }
}
\end{lstlisting}

\subsubsection{3.3
速度测量插件}\label{ux901fux5ea6ux6d4bux91cfux63d2ux4ef6}

使用speed-measure-webpack-plugin测量Webpack构建速度。

\begin{lstlisting}
// vue.config.js
const SpeedMeasurePlugin = require("speed-measure-webpack-plugin");
const smp = new SpeedMeasurePlugin();

module.exports = {
  configureWebpack: config => {
    if (process.env.MEASURE) {
      return smp.wrap(config);
    }
  }
};
\end{lstlisting}

\subsubsection{3.4
压缩与优化插件}\label{ux538bux7f29ux4e0eux4f18ux5316ux63d2ux4ef6}

优化生产环境包体积的插件。

\begin{lstlisting}
// vue.config.js
const CompressionPlugin = require('compression-webpack-plugin');
const TerserPlugin = require('terser-webpack-plugin');

module.exports = {
  configureWebpack: config => {
    if (process.env.NODE_ENV === 'production') {
      return {
        plugins: [
          // Gzip压缩
          new CompressionPlugin({
            algorithm: 'gzip',
            test: /\.(js|css|json|txt|html|ico|svg)(\?.*)?$/i,
            threshold: 10240,
            minRatio: 0.8
          })
        ],
        optimization: {
          minimizer: [
            // JS压缩
            new TerserPlugin({
              terserOptions: {
                compress: {
                  arrows: false,
                  collapse_vars: false,
                  comparisons: false,
                  computed_props: false,
                  hoist_props: false,
                  inline: false,
                  loops: false,
                  negate_iife: false,
                  properties: false,
                  reduce_funcs: false,
                  reduce_vars: false,
                  switches: false,
                  typeofs: false,
                  drop_console: true,
                  drop_debugger: true,
                  pure_funcs: ['console.log']
                },
                mangle: {
                  safari10: true
                }
              },
              parallel: true,
              extractComments: false
            })
          ]
        }
      };
    }
  }
};
\end{lstlisting}

\subsection{4.
性能优化最佳实践}\label{ux6027ux80fdux4f18ux5316ux6700ux4f73ux5b9eux8df5}

\subsubsection{4.1
资源加载优化}\label{ux8d44ux6e90ux52a0ux8f7dux4f18ux5316}

加速智慧水利平台资源加载的关键策略。

\paragraph{资源预加载与预连接}\label{ux8d44ux6e90ux9884ux52a0ux8f7dux4e0eux9884ux8fdeux63a5}

\begin{lstlisting}[language=HTML]
<!-- index.html -->
<head>
  <!-- 预连接到API域名 -->
  <link rel="preconnect" href="https://api.water-platform.com">
  
  <!-- 预加载关键CSS -->
  <link rel="preload" href="/css/critical.css" as="style">
  
  <!-- 预加载字体 -->
  <link rel="preload" href="/fonts/custom-icon.woff2" as="font" type="font/woff2" crossorigin>
  
  <!-- 预取可能需要的资源 -->
  <link rel="prefetch" href="/js/water-monitoring.chunk.js">
</head>
\end{lstlisting}

\paragraph{图像优化策略}\label{ux56feux50cfux4f18ux5316ux7b56ux7565}

\begin{lstlisting}[language=HTML]
<!-- 响应式图像 -->
<img 
  src="/images/reservoir-small.jpg"
  srcset="/images/reservoir-small.jpg 600w,
          /images/reservoir-medium.jpg 1200w,
          /images/reservoir-large.jpg 2000w"
  sizes="(max-width: 600px) 100vw,
         (max-width: 1200px) 50vw,
         33vw"
  alt="水库照片">

<!-- 对现代浏览器使用WebP格式 -->
<picture>
  <source srcset="/images/map.webp" type="image/webp">
  <source srcset="/images/map.jpg" type="image/jpeg">
  <img src="/images/map.jpg" alt="流域地图">
</picture>

<!-- 懒加载非首屏图像 -->
<img 
  loading="lazy"
  src="/images/chart-preview.jpg"
  alt="历史数据图表">
\end{lstlisting}

\paragraph{JavaScript优化}\label{javascriptux4f18ux5316}

\begin{enumerate}
\def\labelenumi{\arabic{enumi}.}
\item
  \textbf{异步加载非关键脚本}

\begin{lstlisting}[language=HTML]
<!-- 异步加载分析脚本 -->
<script async src="/js/analytics.js"></script>

<!-- 延迟加载非关键脚本 -->
<script defer src="/js/non-critical.js"></script>
\end{lstlisting}
\item
  \textbf{动态导入}

\begin{lstlisting}
// 按需加载模块
async function loadMapWhenNeeded() {
  const { initMap } = await import(/* webpackChunkName: "map" */ './map.js');
  initMap();
}

// 按条件懒加载组件
const AdminPanel = () => process.env.NODE_ENV === 'development' 
  ? import('./AdminPanel.vue')
  : import('./ProductionAdminPanel.vue');
\end{lstlisting}
\end{enumerate}

\subsubsection{4.2
运行时性能优化}\label{ux8fd0ux884cux65f6ux6027ux80fdux4f18ux5316}

优化智慧水利平台前端应用的运行时性能。

\paragraph{避免内存泄漏}\label{ux907fux514dux5185ux5b58ux6cc4ux6f0f}

\begin{lstlisting}
// 组件中管理事件监听器
export default {
  data() {
    return {
      waterLevelSocket: null,
      resizeObserver: null
    };
  },
  
  mounted() {
    // WebSocket连接
    this.waterLevelSocket = new WebSocket('wss://api.example.com/water-data');
    this.waterLevelSocket.addEventListener('message', this.handleWaterData);
    
    // ResizeObserver
    this.resizeObserver = new ResizeObserver(this.handleResize);
    this.resizeObserver.observe(this.$refs.chart);
    
    // 窗口事件
    window.addEventListener('online', this.handleOnline);
  },
  
  beforeDestroy() {
    // 清理WebSocket
    if (this.waterLevelSocket) {
      this.waterLevelSocket.removeEventListener('message', this.handleWaterData);
      this.waterLevelSocket.close();
    }
    
    // 清理ResizeObserver
    if (this.resizeObserver) {
      this.resizeObserver.disconnect();
    }
    
    // 清理窗口事件
    window.removeEventListener('online', this.handleOnline);
  }
};
\end{lstlisting}

\paragraph{虚拟滚动大数据列表}\label{ux865aux62dfux6edaux52a8ux5927ux6570ux636eux5217ux8868}

\begin{lstlisting}
<!-- WaterStationList.vue -->
<template>
  <div class="station-list-container">
    <RecycleScroller
      class="scroller"
      :items="stations"
      :item-size="60"
      key-field="id"
      v-slot="{ item }"
    >
      <div class="station-item">
        <div class="station-name">{{ item.name }}</div>
        <div class="water-level" :class="getLevelClass(item)">
          {{ item.waterLevel }}m
        </div>
        <div class="timestamp">{{ formatTime(item.timestamp) }}</div>
      </div>
    </RecycleScroller>
  </div>
</template>

<script>
import { RecycleScroller } from 'vue-virtual-scroller';
import 'vue-virtual-scroller/dist/vue-virtual-scroller.css';

export default {
  components: {
    RecycleScroller
  },
  
  props: {
    stations: {
      type: Array,
      required: true
    }
  },
  
  methods: {
    getLevelClass(station) {
      if (station.waterLevel >= station.warningLevel) {
        return 'warning';
      }
      return 'normal';
    },
    
    formatTime(timestamp) {
      return new Date(timestamp).toLocaleTimeString();
    }
  }
};
</script>
\end{lstlisting}

\paragraph{避免频繁DOM更新}\label{ux907fux514dux9891ux7e41domux66f4ux65b0}

\begin{lstlisting}
<!-- RainfallChart.vue -->
<template>
  <div class="rainfall-chart">
    <div ref="chartContainer" class="chart-container"></div>
    <div class="chart-controls">
      <button @click="debounceRefresh">刷新数据</button>
      <select v-model="timeRange" @change="debounceRangeChange">
        <option value="1h">1小时</option>
        <option value="24h">24小时</option>
        <option value="7d">7天</option>
      </select>
    </div>
  </div>
</template>

<script>
import * as echarts from 'echarts/core';
import _ from 'lodash';

export default {
  data() {
    return {
      chart: null,
      timeRange: '24h',
      debounceRefresh: null,
      debounceRangeChange: null
    };
  },
  
  created() {
    // 防抖处理
    this.debounceRefresh = _.debounce(this.refreshData, 300);
    this.debounceRangeChange = _.debounce(this.handleRangeChange, 300);
  },
  
  mounted() {
    // 初始化图表
    this.chart = echarts.init(this.$refs.chartContainer);
    this.refreshData();
    
    // 使用ResizeObserver替代window事件
    const resizeObserver = new ResizeObserver(_.throttle(() => {
      this.chart.resize();
    }, 100));
    
    resizeObserver.observe(this.$refs.chartContainer);
    this.resizeObserver = resizeObserver;
  },
  
  methods: {
    async refreshData() {
      // 显示加载状态
      this.chart.showLoading();
      
      try {
        // 批量获取数据
        const data = await this.fetchRainfallData(this.timeRange);
        
        // 一次性更新视图
        this.chart.setOption({
          // 图表配置
        });
      } finally {
        this.chart.hideLoading();
      }
    },
    
    handleRangeChange() {
      this.refreshData();
    }
  },
  
  beforeDestroy() {
    if (this.chart) {
      this.chart.dispose();
    }
    
    if (this.resizeObserver) {
      this.resizeObserver.disconnect();
    }
  }
};
</script>
\end{lstlisting}

\subsubsection{4.3
网络性能优化}\label{ux7f51ux7edcux6027ux80fdux4f18ux5316}

优化智慧水利平台的数据传输效率。

\paragraph{API请求优化}\label{apiux8bf7ux6c42ux4f18ux5316}

\begin{lstlisting}
// api-client.js
import axios from 'axios';

class ApiClient {
  constructor() {
    this.client = axios.create({
      baseURL: process.env.VUE_APP_API_BASE_URL,
      timeout: 15000
    });
    
    this.setupInterceptors();
    this.cache = new Map();
  }
  
  setupInterceptors() {
    // 请求拦截器
    this.client.interceptors.request.use(config => {
      // 添加取消令牌
      const source = axios.CancelToken.source();
      config.cancelToken = source.token;
      
      // 存储取消函数
      const requestId = this.getRequestId(config);
      this.abortControllers = this.abortControllers || {};
      
      // 如果存在相同请求，取消旧请求
      if (this.abortControllers[requestId]) {
        this.abortControllers[requestId]();
      }
      
      this.abortControllers[requestId] = source.cancel;
      
      return config;
    });
    
    // 响应拦截器
    this.client.interceptors.response.use(
      response => {
        // 清理取消函数
        const requestId = this.getRequestId(response.config);
        delete this.abortControllers[requestId];
        
        return response;
      },
      error => {
        // 如果是取消的请求，不传播错误
        if (axios.isCancel(error)) {
          return new Promise(() => {});
        }
        
        return Promise.reject(error);
      }
    );
  }
  
  getRequestId(config) {
    return `${config.method}:${config.url}:${JSON.stringify(config.params)}`;
  }
  
  // 带缓存的GET请求
  async getCached(url, params, cacheTime = 60000) {
    const cacheKey = `${url}:${JSON.stringify(params)}`;
    const cached = this.cache.get(cacheKey);
    
    if (cached && Date.now() - cached.timestamp < cacheTime) {
      return cached.data;
    }
    
    const response = await this.client.get(url, { params });
    
    this.cache.set(cacheKey, {
      data: response.data,
      timestamp: Date.now()
    });
    
    return response.data;
  }
  
  // 批量请求
  async batchGet(requests) {
    return Promise.all(
      requests.map(req => this.client.get(req.url, { params: req.params }))
    );
  }
}

export default new ApiClient();
\end{lstlisting}

\paragraph{数据压缩}\label{ux6570ux636eux538bux7f29}

\begin{lstlisting}
// 服务器端压缩配置 (Express示例)
const compression = require('compression');
app.use(compression());

// 前端检测压缩支持
function checkCompressionSupport() {
  const headers = new Headers();
  headers.append('Accept-Encoding', 'gzip, deflate, br');
  
  return fetch('/api/check-compression', { headers })
    .then(response => {
      const encoding = response.headers.get('Content-Encoding');
      return !!encoding;
    })
    .catch(() => false);
}
\end{lstlisting}

\subsection{总结}\label{ux603bux7ed3}

前端性能优化工具是智慧水利平台提升用户体验的关键要素。通过浏览器开发者工具进行性能分析，使用监控和错误追踪工具实时掌握应用状态，利用打包分析工具优化资源加载，可以全面提升智慧水利平台的性能表现。

在实际开发中，应根据平台的实际需求和使用场景，选择合适的工具和优化策略。特别是对于数据可视化和实时监控功能，要重点关注资源加载优化、运行时性能和网络传输效率，确保即使在网络条件较差的场景下也能提供良好的用户体验。通过持续的性能监控和优化，智慧水利平台可以为用户提供高效、流畅的操作体验。
